\section{Data}
\label{sec:data}

\subsection{Source corpus}
All training data originate from the DeepMath-103K
collection~\cite{deepmath103k}, distributed in our project as
\code{data/raw/deepmath\_103k\_raw.jsonl}. Each record contains a problem
statement, a final reference answer, up to three candidate ``R1-style''
solution traces, a difficulty score, and a topic path. A minimal example is
shown below.

\begin{lstlisting}[language=json,caption={A raw record from
\code{deepmath\_103k\_raw.jsonl}.}]
{
  "question": "Evaluate the limit: ...",
  "final_answer": "0",
  "r1_solution_1": "Okay, so I have this limit to evaluate...",
  "r1_solution_2": "...",
  "r1_solution_3": "...",
  "difficulty": 4.5,
  "topic": "Mathematics -> ..."
}
\end{lstlisting}

Two distinct downstream artefacts are extracted from this source: a small
SFT set that requires a clean reasoning trace, and a much larger RL-ready
set that requires only a verifiable final answer.

\subsection{SFT data: \code{strict\_main}}
The SFT subset is created by keeping only records whose problem,
final answer, and at least one candidate trace all pass strict cleanliness
filters. The resulting set, \code{strict\_main}, contains \textbf{1\,298}
training samples and \textbf{169} validation samples. Each sample is then
exported in a ShareGPT-style \code{messages} format, where the user message
contains the math problem under a unified instruction template, and the
assistant message contains a cleaned reasoning trace ending in a single
\code{Final Answer:} line. Trace cleaning includes (i) removing
\code{</think>} markers, (ii) dropping repeated \code{Final Answer} headings
or duplicated answer lines, (iii) stripping unnecessary \code{\textbackslash
boxed\{\}} wrappers, and (iv) appending a single normalised
\code{Final Answer: <answer>} line.

\subsection{RL-ready data}
The RL data is generated by \code{data/build\_rl\_ready\_dataset.py}, which
relaxes the SFT requirement of ``must have a complete reasoning trace'' to
``the question is well-formed and the final answer can be verified''. Three
variants are kept:

\begin{table}[h]
\centering
\small
\begin{tabular}{lrrl}
\toprule
\textbf{Name} & \textbf{Train} & \textbf{Valid} & \textbf{Notes} \\
\midrule
\code{rl\_strict}            & 81\,463  & 10\,207 & strict length limits, main RL pool \\
\code{rl\_relaxed}           & 81\,478  & 10\,208 & relaxed length limits, ablation pool \\
\code{bigmath\_verified\_grpo} & 194\,274 & 24\,489 & alternative larger pool, not used here \\
\bottomrule
\end{tabular}
\caption{The three RL-ready variants. This report uses
\code{rl\_strict} and a balanced sub-sample of it.}
\label{tab:rl-data}
\end{table}

\subsection{The 24K balanced sub-sample}
To control training time on a single GPU we additionally derive a
\textbf{24K balanced sub-sample} from \code{rl\_strict}, called
\code{sampled24k}. The sampling has two stages:

\begin{enumerate}
  \item \emph{Difficulty buckets.} Examples are split into four buckets by
    the \code{difficulty} field: \code{low}~$\leq 3.0$, \code{mid}~$\leq
    6.0$, \code{high}~$\leq 8.0$, and \code{very\_high}~$> 8.0$ (or
    missing). Each bucket gets a fixed quota: 2\,400 / 10\,800 / 7\,200 /
    3\,600.
  \item \emph{Task-type stratification within bucket.} Inside each bucket,
    the per-task-type quota follows the original task-type proportion of
    that bucket, with leftover quota redistributed proportionally if a
    task type is short of samples.
\end{enumerate}

The validation split is the entire \code{rl\_strict\_valid}
(10\,207 examples). A fixed random seed (\code{seed=42}) makes the sample
reproducible. The resulting 24K mixture is dominated by \code{algebra} and
\code{calculus}, with non-trivial coverage of \code{geometry},
\code{number\_theory}, and \code{probability}.

\subsection{Data layers}
For clarity it is useful to separate the data into four layers, summarised
in \cref{tab:layers}.

\begin{table}[h]
\centering
\small
\begin{tabularx}{\linewidth}{lXl}
\toprule
\textbf{Layer} & \textbf{Purpose} & \textbf{Example file} \\
\midrule
1. Raw download   & Original DeepMath records              & \code{data/raw/deepmath\_103k\_raw.jsonl} \\
2. Cleaned middle & Normalised question, answer, candidates & \code{data/processed/filtered\_strict\_train.jsonl} \\
3. SFT export     & ShareGPT \code{messages} format         & \code{data/processed/sft\_strict\_main\_train.json} \\
4. RL export      & Single prompt $+$ verifiable target answer & \code{data/processed/grpo\_strict\_main\_train.jsonl} \\
\bottomrule
\end{tabularx}
\caption{Four data layers used in this project.}
\label{tab:layers}
\end{table}

\begin{remark}
The clean separation between layer 3 and layer 4 is what makes the SFT and
RL stages compatible. Both stages use the same instruction template and
the same final-answer extraction logic, so a model trained at layer 3 can
be directly continued with RL at layer 4.
\end{remark}
